\scanpage{9}
et si
\[
\eta_p\in H^1(Q_p,G/H)
\]
est la classe de ce torseur, alors par \unsure{du general} nonsense on trouve que
\[
\xi_p\ (\text{image de }\xi\text{ dans }H^2(Q_p,H))=\partial(\eta_p)
\]
où \deleted{$\unsure{\eta_p}$} $\partial$ \unsure{est la} suite ex. relative à :
\diagram{\begin{tikzcd}
0\arrow[r]&H\arrow[r,dash]&G\arrow[r,dash]&G/H\arrow[r]&0
\end{tikzcd}}
Par suite la classe $\xi_p$ est déterminée entièrement par la connaissance de $\eta_p$. Or :
\diagram{\begin{tikzcd}
\unsure{0}\arrow[r]&G_m\arrow[r,dash]\arrow[d,equal] &G/H\arrow[r,dash]&\widehat{\mathbf Z}\arrow[r,dash]&\unsure{0}\\
&(G/H)^0&&&
\end{tikzcd}}
(extension qui splitte)

donc
\[
H^1(Q_p,G/H)\xrightarrow{\sim}H^1(Q_p,\widehat{\mathbf Z})
\]
donc il suffit de décrire l'image de $\eta_p$ dans $H^1(Q_p,\widehat{\mathbf Z})$, i.e. encore le torseur des isom de $(F_0)_{Q_p}\mid\operatorname{Niv}^0(\mathcal M)^0$ avec $T_p\mid\operatorname{Niv}^0(\mathcal M)^0$. On trouve

(histoire générale valable sur tt corps parfait $k$, en travaillant avec \deleted{$\unsure{F}$} $T_p$ sur \deleted{\unsure{le}} $\operatorname{frac}(W(k))$...)

l'élément de $H^1(Q_p,\widehat{\mathbf Z})$ \inserted{$\simeq\operatorname{Hom}(\pi_1(Q_p),\widehat{\mathbf Z})$} correspondant à :

l'isom $\operatorname{Gal}(\widetilde Q_p/Q_p)\simeq\widehat{\mathbf Z}$ donné par Frobenius.

Notons que les composantes de $\xi$ suivant $H^2(Q,H^0)$ et $H^2(Q,\mu_2)$ sont connues quand on connaît leurs composantes \uline{en ttes places} (y compris l'$\infty$), et pour la dernière, à cause de la formule $\sum\operatorname{inv}_p=0$, il suffit de la connaître \unsure{à} distance finie.
\margin{\unsure{vif|vérif}}
\mspage{8}
